% GENERATED FILE. Edit canonical lessons, then run npm run generate:books.
% canonical-source-hash: fnv1a64:75b6915c31d2d27a
% canonical-lessons: JA-C196-mukau, JA-C196-kinjo, JA-C196-juusho, JA-C196-chizu, JA-C196-mawaru

\chapter{Addresses, Maps and Neighbours}
\label{ch:ja-addresses-maps-and-neighbours}
\hlchaptermodality{196}

\begin{chapteropening}
\textbf{By the end of this chapter:} \emph{I can talk about my neighbourhood and find an address.}

Complete the last lesson of chapter 196: I can talk about my neighbourhood and find an address.
\end{chapteropening}

\section[\texorpdfstring{\ja{むかう} (mukau)}{むかう (mukau)}]{\ja{むかう} (mukau) — to head for; to face}

\label{lesson:JA-C196-mukau}



\begin{quote}
Before the new one: say the Japanese for to go past, then the Japanese for to go back.
\end{quote}

\subsection*{You'll want to know: \ja{むかう}}
\textbf{\ja{むかう}} — \emph{mukau} — \textquotedblleft{}to head for; to face\textquotedblright{}.

\begin{quote}
\textbf{\ja{いま} \ja{えきに} \ja{むかいます。}} — \emph{ima eki ni mukaimasu} — \textquotedblleft{}I am heading for the station now.\textquotedblright{}
\end{quote}

Polite \textbf{\ja{むかいます}}, \emph{mukaimasu}; a godan verb.

\begin{cousinweb}
\textbf{\ja{むかう}} (\emph{mukau}) is a native Japanese word (wago), inherited rather than borrowed from Chinese; its kanji was attached later as a kun'yomi (JMdict, KANJIDIC2).
\end{cousinweb}

\begin{grammarlens}[title={What you've built}]
One more word for everyday life at A2.
\end{grammarlens}

\subsection*{Your turn}
\begin{itemize}
\raggedright
  \item \emph{Say it:} \emph{mukau}
  \item \emph{Say it:} \emph{mukau}, once more
  \item \emph{Say it:} \emph{ima eki ni mukaimasu} — I am heading for the station now.
  \item \emph{Recall:} say the Japanese for to go past, then the Japanese for to go back, then say \emph{mukau} again
\end{itemize}

\begin{tcolorbox}[breakable,colback=teal!4,colframe=teal!35!black,title={Before you move on}]
What does \ja{むかう} mean? (\textquotedblleft{}To head for; to face\textquotedblright{}.) Say it once more.
\end{tcolorbox}
\section[\texorpdfstring{\ja{きんじょ} (kinjo)}{きんじょ (kinjo)}]{\ja{きんじょ} (kinjo) — neighbourhood; vicinity}

\label{lesson:JA-C196-kinjo}



\begin{quote}
Before the new one: say the Japanese for to go back, then the Japanese for to head for.
\end{quote}

\subsection*{You'll want to know: \ja{きんじょ}}
\textbf{\ja{きんじょ}} — \emph{kinjo} — \textquotedblleft{}neighbourhood; vicinity\textquotedblright{}.

\begin{quote}
\textbf{\ja{きんじょに} \ja{いい} \ja{みせが} \ja{あります。}} — \emph{kinjo ni ii mise ga arimasu} — \textquotedblleft{}There is a good shop in my neighbourhood.\textquotedblright{}
\end{quote}

\begin{cousinweb}
\textbf{\ja{きんじょ}} (\emph{kinjo}) is Sino-Japanese (kango): two Chinese characters read with their on'yomi, readings borrowed from Chinese. They are \emph{kin}, \textquotedblleft{}near\textquotedblright{} (Middle Chinese \emph{gj+nX}) and \emph{jo}, \textquotedblleft{}place\textquotedblright{} (\emph{srjoX}) (KANJIDIC2; Middle Chinese in Baxter's transcription).
\end{cousinweb}

\begin{grammarlens}[title={What you've built}]
One more word for everyday life at A2.
\end{grammarlens}

\subsection*{Your turn}
\begin{itemize}
\raggedright
  \item \emph{Say it:} \emph{kinjo}
  \item \emph{Say it:} \emph{kinjo}, once more
  \item \emph{Say it:} \emph{kinjo ni ii mise ga arimasu} — There is a good shop in my neighbourhood.
  \item \emph{Recall:} say the Japanese for to go back, then the Japanese for to head for, then say \emph{kinjo} again
\end{itemize}

\begin{tcolorbox}[breakable,colback=teal!4,colframe=teal!35!black,title={Before you move on}]
What does \ja{きんじょ} mean? (\textquotedblleft{}Neighbourhood; vicinity\textquotedblright{}.) Say it once more.
\end{tcolorbox}
\section[\texorpdfstring{\ja{じゅうしょ} (jūsho)}{じゅうしょ (jūsho)}]{\ja{じゅうしょ} (jūsho) — address; residence}

\label{lesson:JA-C196-juusho}



\begin{quote}
Before the new one: say the Japanese for to head for, then the Japanese for neighbourhood.
\end{quote}

\subsection*{You'll want to know: \ja{じゅうしょ}}
\textbf{\ja{じゅうしょ}} — \emph{jūsho} — \textquotedblleft{}address; residence\textquotedblright{}.

\begin{quote}
\textbf{\ja{てがみに} \ja{じゅうしょを} \ja{かきます。}} — \emph{tegami ni jūsho o kakimasu} — \textquotedblleft{}I write the address on the letter.\textquotedblright{}
\end{quote}

\begin{cousinweb}
\textbf{\ja{じゅうしょ}} (\emph{jūsho}) is Sino-Japanese (kango): two Chinese characters read with their on'yomi, readings borrowed from Chinese. They are \emph{jū}, \textquotedblleft{}dwell\textquotedblright{} (Middle Chinese \emph{drjuH}) and \emph{sho}, \textquotedblleft{}place\textquotedblright{} (\emph{srjoX}) (KANJIDIC2; Middle Chinese in Baxter's transcription).
\end{cousinweb}

\begin{grammarlens}[title={What you've built}]
One more word for everyday life at A2.
\end{grammarlens}

\subsection*{Your turn}
\begin{itemize}
\raggedright
  \item \emph{Say it:} \emph{jūsho}
  \item \emph{Say it:} \emph{jūsho}, once more
  \item \emph{Say it:} \emph{tegami ni jūsho o kakimasu} — I write the address on the letter.
  \item \emph{Recall:} say the Japanese for to head for, then the Japanese for neighbourhood, then say \emph{jūsho} again
\end{itemize}

\begin{tcolorbox}[breakable,colback=teal!4,colframe=teal!35!black,title={Before you move on}]
What does \ja{じゅうしょ} mean? (\textquotedblleft{}Address; residence\textquotedblright{}.) Say it once more.
\end{tcolorbox}
\section[\texorpdfstring{\ja{ちず} (chizu)}{ちず (chizu)}]{\ja{ちず} (chizu) — map}

\label{lesson:JA-C196-chizu}



\begin{quote}
Before the new one: say the Japanese for neighbourhood, then the Japanese for address.
\end{quote}

\subsection*{You'll want to know: \ja{ちず}}
\textbf{\ja{ちず}} — \emph{chizu} — \textquotedblleft{}map\textquotedblright{}.

\begin{quote}
\textbf{\ja{ちずで} \ja{みちを} \ja{さがします。}} — \emph{chizu de michi o sagashimasu} — \textquotedblleft{}I look for the way on a map.\textquotedblright{}
\end{quote}

\begin{cousinweb}
\textbf{\ja{ちず}} (\emph{chizu}) is Sino-Japanese (kango): two Chinese characters read with their on'yomi, readings borrowed from Chinese. They are \emph{chi}, \textquotedblleft{}ground, earth\textquotedblright{} (Middle Chinese \emph{dijH}) and \emph{zu}, \textquotedblleft{}drawing, plan\textquotedblright{} (\emph{du}) (KANJIDIC2; Middle Chinese in Baxter's transcription).
\end{cousinweb}

\begin{grammarlens}[title={What you've built}]
One more word for everyday life at A2.
\end{grammarlens}

\subsection*{Your turn}
\begin{itemize}
\raggedright
  \item \emph{Say it:} \emph{chizu}
  \item \emph{Say it:} \emph{chizu}, once more
  \item \emph{Say it:} \emph{chizu de michi o sagashimasu} — I look for the way on a map.
  \item \emph{Recall:} say the Japanese for neighbourhood, then the Japanese for address, then say \emph{chizu} again
\end{itemize}

\begin{tcolorbox}[breakable,colback=teal!4,colframe=teal!35!black,title={Before you move on}]
What does \ja{ちず} mean? (\textquotedblleft{}Map\textquotedblright{}.) Say it once more.
\end{tcolorbox}
\section[\texorpdfstring{\ja{まわる} (mawaru)}{まわる (mawaru)}]{\ja{まわる} (mawaru) — to go round; to make the rounds, to spin}

\label{lesson:JA-C196-mawaru}



\begin{quote}
Before the new one: say the Japanese for address, then the Japanese for map.
\end{quote}

\subsection*{You'll want to know: \ja{まわる}}
\textbf{\ja{まわる}} — \emph{mawaru} — \textquotedblleft{}to go round; to make the rounds, to spin\textquotedblright{}.

\begin{quote}
\textbf{\ja{しょうてんがいの} \ja{みせを} \ja{まわります。}} — \emph{shōtengai no mise o mawarimasu} — \textquotedblleft{}I go round the shops on the shopping street.\textquotedblright{}
\end{quote}

Polite \textbf{\ja{まわります}}, \emph{mawarimasu}; a godan verb. Its transitive partner \textbf{\ja{まわす}} is taught already.

\begin{cousinweb}
\textbf{\ja{まわる}} (\emph{mawaru}) is a native Japanese word (wago), inherited rather than borrowed from Chinese; its kanji was attached later as a kun'yomi (JMdict, KANJIDIC2).
\end{cousinweb}

\begin{grammarlens}[title={What you've built}]
One more word for everyday life at A2.
\end{grammarlens}

\subsection*{Your turn}
\begin{itemize}
\raggedright
  \item \emph{Say it:} \emph{mawaru}
  \item \emph{Say it:} \emph{mawaru}, once more
  \item \emph{Say it:} \emph{shōtengai no mise o mawarimasu} — I go round the shops on the shopping street.
  \item \emph{Recall:} say the Japanese for address, then the Japanese for map, then say \emph{mawaru} again
\end{itemize}

\begin{tcolorbox}[breakable,colback=teal!4,colframe=teal!35!black,title={Before you move on}]
What does \ja{まわる} mean? (\textquotedblleft{}To go round; to make the rounds, to spin\textquotedblright{}.) Say it once more.
\end{tcolorbox}
